% 09. kapitulua - Ariketak: ekosistema agenteekin
% Helburua: jatorrizko taldea (Nami, Ane, Mikel, Gotzon), orain agenteekin. Gotzonen gainbegiratze-lana agente batek hartzen du;
%           Nami eta Mikel kodetzeko agenteei eskuordetzen diete; Ane giza produktu-jabea da.
% Aurreko liburutik berrerabil daitekeena:
%   - aurreko_liburua_2024/kapituluak/9-ariketa_adarkatze.tex (narrazio osoa)
%   - adibideak/ekosistema_maven (Java/Maven proiektua, test-karpetarekin), adibideak/ekosistema_java, adibideak/ekosistema_python
% Irudiak: Pictures/09/ (aurreko 9. kapituluko irudi guztiak: zirriborroak, backlog, evolution, github-flow, bugfix, quickfix, integrazioa; goiburu hutsa header_09.pdf.xoj)
\todo[inline]{Laburpena (behin-behinekoa): Kapitulu honetan proiektu oso bat garatuko dugu Scrumban agentikoarekin, sprint batzuetan zehar: aurreko eskuliburuko ekosistemaren simulazioa, orain lantalde hibrido batekin. Gizakiek \textit{product backlog}-a idatziko dute eta sprinta planifikatuko dute; agenteek zereginak hartuko dituzte eta \textit{pull request}-ak irekiko dituzte; scrumban-agenteak fluxua zainduko du eta blokeoak eskalatuko ditu; gizakiek berrikusi, bateratu eta sprint-berrikuspena egingo dute. Ariketa bakoitzak egoera erreal bat planteatzen du ---agente bik fitxategi bera ukitu dute, zeregin bat blokeatuta geratu da, berrikuspenen ilara luzeegia da--- eta taldeak erabaki behar du zer egin. Giten historia-grafoa eta proiektu-taula marraztuko ditugu une bakoitzean, aurreko eskuliburuan egin genuen bezala.}

\section{Ariketa ebatziak}

\section{Proposatutako ariketa gehigarriak}
